\section{Машинный способ задания функций. Коммутативные диаграммы}

\tikzset{blk/.style={draw,thick,minimum width=7mm,minimum height=7mm,inner sep=1pt},
         wire/.style={-{Latex[length=1.6mm]}},
         inp/.style={circle,fill,inner sep=0.9pt},
         schema/.style={draw,dashed,inner sep=4mm}}

Аналитический способ задания функций, формализованный в предыдущем
параграфе с помощью понятий суперпозиции и $H$-реализуемой функции,
предполагает наличие некоторой системы функций $H$. Функции, входящие в $H$,
также могут оказаться суперпозициями каких-то других функций, а те, в свою
очередь, выражаться через еще более простые функции и~т.\,д. Однако в любом
случае придем к некоторой <<простейшей>> системе $H^0$, которую надо будет
задавать уже с помощью каких-то иных средств. Например, для множества $C^R$
целых рациональных функций простейшей системой является
\begin{equation}\tag{1.16}
H_1=\{x_1+x_2;\ x_1x_2;\ a,\ \forall a\in\R\}.
\end{equation}

\begin{figure}[h]
\centering
\begin{tikzpicture}[x=1cm,y=1cm]
 \node[blk] (p) at (0,0) {$+$};
 \draw[wire] (-1,0.2) -- (p.west|-0,0.2); \draw[wire] (-1,-0.2) -- (p.west|-0,-0.2);
 \node[blk] (m) at (2.3,0) {$\times$};
 \draw[wire] (p.east) -- (1.25,0);
 \draw[wire] (1.3,0.2) -- (m.west|-0,0.2); \draw[wire] (1.3,-0.2) -- (m.west|-0,-0.2);
 \draw[wire] (m.east) -- (3.3,0);
 \node[blk] (a) at (1.15,-1.3) {$a$};
 \draw[wire] (a.east) -- (2.1,-1.3);
\end{tikzpicture}
\figcap{Рис. 9.}
\end{figure}

Символы операций, входящих в эту систему, могут быть поняты потому, что мы
помним таблицы сложения и умножения, и знаем правила выполнения этих операций
для любых $x$ и $y$. Поэтому можно считать, что эти операции некоторым
образом материализованы в нашем мозгу. Впрочем, можно представить их в виде
блоков некоторых аналоговых устройств (<<функциональных преобразователей>>),
блоков ЭВМ, подпрограмм, таблиц, номограмм и~т.\,д. Эти блоки будем условно
\bookpage{30}%
изображать так, как показано на рис.~9. Такой способ задания функций,
независимо от способа его реализации, будем называть \emph{машинным}. Из
сказанного следует, что аналитическое задание функций возможно лишь при
наличии некоторой системы функций, заданных машинным способом.

Используя изображение базисных функций в виде блоков, можно дать простую
интерпретацию операции суперпозиции функций. Пусть базисная система состоит
из трех функций
\[
H=\{\varphi(x,y),\ \psi(x,y),\ f(x,y,z)\},
\]
изображенных в виде блоков (рис.~10). Тогда, например, суперпозициям
\[
w_1=f[x,\ \varphi(x,y),\ \psi(z,t)];\qquad
w_2=\varphi[f(x,x,z),\ \varphi(y,z)]
\]
будут соответствовать две функциональные схемы, показанные на рис.~11.

\begin{figure}[h]
\centering
\begin{tikzpicture}[x=1cm,y=1cm]
 \node[blk,minimum size=9mm] (a) at (0,0) {$\varphi$};
 \draw[wire] (-1,0.2) -- (a.west|-0,0.2); \draw[wire] (-1,-0.2) -- (a.west|-0,-0.2);
 \draw[wire] (a.east) -- (0.95,0);
 \node[blk,minimum size=9mm] (b) at (2.2,0) {$\psi$};
 \draw[wire] (1.2,0.2) -- (b.west|-0,0.2); \draw[wire] (1.2,-0.2) -- (b.west|-0,-0.2);
 \draw[wire] (b.east) -- (3.2,0);
 \node[blk,minimum size=9mm] (c) at (1.1,-1.5) {$f$};
 \draw[wire] (0.1,-1.25) -- (c.west|-0,-1.25); \draw[wire] (0.1,-1.5) -- (c.west|-0,-1.5);
 \draw[wire] (0.1,-1.75) -- (c.west|-0,-1.75);
 \draw[wire] (c.east) -- (2.1,-1.5);
\end{tikzpicture}
\figcap{Рис. 10.}
\end{figure}

\begin{figure}[h]
\centering
\begin{tikzpicture}[x=1cm,y=1cm]
 % --- а: w1 = f[x, phi(x,y), psi(z,t)]
 \node[inp,label=left:$x$] (x) at (0,1.2) {};
 \node[inp,label=left:$y$] (y) at (0,0.9) {};
 \node[inp,label=left:$z$] (z) at (0,-0.3) {};
 \node[inp,label=left:$t$] (t) at (0,-0.6) {};
 \node[blk] (ph) at (1.6,1.05) {$\varphi$};
 \node[blk] (ps) at (1.6,-0.45) {$\psi$};
 \node[blk,minimum height=11mm] (f) at (3.5,0.45) {$f$};
 \draw[wire] (x) -- (ph.west|-x);
 \draw[wire] (y) -- (ph.west|-y);
 \draw[wire] (z) -- (ps.west|-z);
 \draw[wire] (t) -- (ps.west|-t);
 \fill (0.6,1.2) circle (0.9pt);
 \draw[wire] (0.6,1.2) |- (2.6,1.75) |- (f.west|-0,0.75);
 \draw[wire] (ph.east) -| (2.8,0.45) -- (f.west|-0,0.45);
 \draw[wire] (ps.east) -| (2.8,0.15) -- (f.west|-0,0.15);
 \draw[wire] (f.east) -- (4.6,0.45) node[right] {$w_1$};
 \draw[dashed] (0.35,-1.0) rectangle (4.35,2.05);
 \node at (2.35,-1.35) {\emph{а}};
\end{tikzpicture}\qquad
\begin{tikzpicture}[x=1cm,y=1cm]
 % --- б: w2 = phi[f(x,x,z), phi(y,z)]
 \node[inp,label=left:$x$] (x) at (0,1.1) {};
 \node[inp,label=left:$y$] (y) at (0,-0.2) {};
 \node[inp,label=left:$z$] (z) at (0,-0.5) {};
 \node[blk,minimum height=10mm] (f) at (1.5,1.1) {$f$};
 \node[blk] (ph1) at (3.0,-0.35) {$\varphi$};
 \node[blk] (ph2) at (4.4,1.0) {$\varphi$};
 \fill (0.6,1.1) circle (0.9pt);
 \draw[wire] (x) -- (0.6,1.1) |- (f.west|-0,1.35);
 \draw[wire] (0.6,1.1) -- (f.west|-0,1.1);
 \fill (0.75,-0.5) circle (0.9pt);
 \draw[wire] (0.75,-0.5) |- (f.west|-0,0.85);
 \draw[wire] (y) -- (ph1.west|-y);
 \draw[wire] (z) -- (ph1.west|-z);
 \draw[wire] (f.east) -- (ph2.west|-0,1.1);
 \draw[wire] (ph1.east) -| (3.9,0.9) -- (ph2.west|-0,0.9);
 \draw[wire] (ph2.east) -- (5.3,1.0) node[right] {$w_2$};
 \draw[dashed] (0.35,-1.0) rectangle (5.1,1.95);
 \node at (2.7,-1.35) {\emph{б}};
\end{tikzpicture}
\figcap{Рис. 11.}
\end{figure}

Заметим, что если в некоторой схеме заменим входящие в нее элементарные
блоки блоками, реализующими какие-то другие операции (с теми же
аргументами, что и заменяемые), то получим функциональные схемы для
формулы, имеющей эквивалентную структуру (см. \S\,4). Например, схема,
показанная на рис.~11,~\emph{а}, при замене
\begin{equation}\tag{1.17}
\begin{pmatrix}f; & \varphi; & \psi\\ x+y+z; & xy; & z^t\end{pmatrix}
\end{equation}
становится схемой функции $w_1=x+xy+z^t$.

Пусть теперь $f[\varphi_1,\dots,\varphi_s]$ есть некоторая формула,
построенная с помощью системы
$H=\{\varphi_i(x_{1i},\dots,x_{n_ii})\}$, $i=1,\dots,s$, а $\Psi_i$ "---
некоторые множества функций вида $\psi_i(x_{1i},\dots,x_{n_ii})$.
Рассмотрим множество всевозможных замен вида
\begin{equation}\tag{1.18}
\begin{pmatrix}\varphi_1, & \dots, & \varphi_s\\ \psi_1, & \dots, & \psi_s\end{pmatrix},
\end{equation}
получаемых путем выбора элементов $(\psi_1,\dots,\psi_s)$ из множества
$\Psi_1\times\dots\times\Psi_s$. Такие замены приведут к образованию
некоторого
\bookpage{31}%
функционального множества $\mathfrak F=\{f[\psi_1,\dots,\psi_s]\}$.
Формулы, определяющие функции из $\mathfrak F$, имеют, очевидно, одинаковую
структуру. Нетрудно заметить, что операция суперпозиции "--- самая дешевая
с точки зрения технической реализации ее: ей соответствует простое
соединение выходов и входов базисных блоков.

Если функции, входящие в систему $H$, обладают некоторой системой свойств,
таких как коммутативность, ассоциативность, дистрибутивность или
каких-либо других, то можно осуществлять тождественные преобразования
$H$-реализуемых функций и получать, вообще, различные формулы для одной
и той же функции. Следовательно, для нее могут оказаться разными
и функциональные схемы. Например, схемы для полинома
$y=x^3+2x^2+3x+4\equiv((x+2)x+3)x+4$ показаны на рис.~12,~\emph{а},
\emph{б}\fixed{в оригинале под схемами рис.~12 нет подписей \emph{а}
и~\emph{б}; верхняя схема "--- \emph{а}, нижняя "--- \emph{б}.} (для
формул $x^3+2x^2+3x+4$, $((x+2)x+3)x+4$
соответственно).\fixed{в оригинале вторая формула напечатана как
$((x+2)+3)x+4$; это выражение равно $x^2+5x+4$, а не данному полиному.}
В то время как для блок-схемы, приведенной на рис.~12,~\emph{а}, понадобилось
четыре блока умножения, для блок-схемы, показанной на рис.~12,~\emph{б},
"--- только два (при равном числе остальных блоков). На этом примере
показана возможность упрощения на практике вычисли\-%
\bookpage{32}%
тельных устройств (например, с целью удешевления, повышения быстродействия
и~т.\,д.) на уровне тождественных преобразований.

\begin{figure}[h]
\centering
\begin{tikzpicture}[x=1cm,y=1cm,blk/.append style={minimum size=6.5mm}]
 % --- а: x^3 + 2x^2 + 3x + 4
 \node[inp,label=left:$x$] (x) at (0,0) {};
 \node[blk] (m1) at (1.2,0) {$\times$};
 \node[blk] (c2) at (1.2,-1.0) {$2$};
 \node[blk] (c3) at (1.2,-1.85) {$3$};
 \node[blk] (c4) at (1.2,-2.9) {$4$};
 \node[blk] (m2) at (2.5,0.15) {$\times$};
 \node[blk] (m3) at (2.5,-0.95) {$\times$};
 \node[blk] (a1) at (3.8,-0.4) {$+$};
 \node[blk] (m4) at (2.5,-2.15) {$\times$};
 \node[blk] (a2) at (3.8,-2.3) {$+$};
 \node[blk] (a3) at (5.0,-1.35) {$+$};
 \fill (0.45,0) circle (0.9pt);
 \draw[wire] (x) -- (0.45,0) |- (m1.west|-0,0.15);
 \draw[wire] (0.45,0) |- (m1.west|-0,-0.15);
 \draw[wire] (0.45,0) |- (1.85,0.75) |- (m2.west|-0,0.3);
 \fill (1.9,0) circle (0.9pt);
 \draw[wire] (m1.east) -- (1.9,0) |- (m2.west|-0,0.0);
 \draw[wire] (1.9,0) |- (m3.west|-0,-0.8);
 \draw[wire] (c2.east) -| (2.0,-1.1) -- (m3.west|-0,-1.1);
 \draw[wire] (m2.east) -| (3.25,-0.25) -- (a1.west|-0,-0.25);
 \draw[wire] (m3.east) -| (3.25,-0.55) -- (a1.west|-0,-0.55);
 \draw[wire] (c3.east) -| (2.0,-2.0) -- (m4.west|-0,-2.0);
 \draw[wire] (0.45,0) |- (m4.west|-0,-2.3);
 \draw[wire] (m4.east) -| (3.25,-2.15) -- (a2.west|-0,-2.15);
 \draw[wire] (c4.east) -| (3.25,-2.45) -- (a2.west|-0,-2.45);
 \draw[wire] (a1.east) -| (4.45,-1.2) -- (a3.west|-0,-1.2);
 \draw[wire] (a2.east) -| (4.45,-1.5) -- (a3.west|-0,-1.5);
 \draw[wire] (a3.east) -- (5.8,-1.35);
 \draw[dashed] (0.25,-3.4) rectangle (5.6,1.05);
 \node at (2.9,-3.75) {\emph{а}};
\end{tikzpicture}\qquad
\begin{tikzpicture}[x=1cm,y=1cm,blk/.append style={minimum size=6.5mm}]
 % --- б: ((x+2)x+3)x+4
 \node[inp,label=left:$x$] (x) at (0,0) {};
 \node[blk] (c2) at (1.0,-0.75) {$2$};
 \node[blk] (c3) at (1.0,-1.75) {$3$};
 \node[blk] (c4) at (1.0,-2.75) {$4$};
 \node[blk] (a1) at (2.0,-0.15) {$+$};
 \node[blk] (m1) at (3.1,-0.15) {$\times$};
 \node[blk] (a2) at (4.0,-1.3) {$+$};
 \node[blk] (m2) at (4.9,-0.15) {$\times$};
 \node[blk] (a3) at (5.8,-1.6) {$+$};
 \fill (0.4,0) circle (0.9pt);
 \draw[wire] (x) -- (a1.west|-0,0);
 \draw[wire] (c2.east) -| (1.5,-0.3) -- (a1.west|-0,-0.3);
 \draw (0.4,0) |- (4.45,0.55);
 \fill (2.6,0.55) circle (0.9pt);
 \draw[wire] (2.6,0.55) |- (m1.west|-0,0.0);
 \draw[wire] (4.45,0.55) |- (m2.west|-0,0.0);
 \draw[wire] (a1.east) -- (2.5,-0.15) |- (m1.west|-0,-0.3);
 \draw[wire] (m1.east) -- (3.55,-0.15) |- (a2.west|-0,-1.15);
 \draw[wire] (c3.east) -| (3.45,-1.45) -- (a2.west|-0,-1.45);
 \draw[wire] (a2.east) -- (4.45,-1.3) |- (m2.west|-0,-0.3);
 \draw[wire] (m2.east) -- (5.38,-0.15) |- (a3.west|-0,-1.45);
 \draw[wire] (c4.east) -| (5.3,-1.75) -- (a3.west|-0,-1.75);
 \draw[wire] (a3.east) -- (6.6,-1.6);
 \draw[dashed] (0.2,-3.2) rectangle (6.4,0.85);
 \node at (3.3,-3.55) {\emph{б}};
\end{tikzpicture}
\figcap{Рис. 12.}
\end{figure}

Сказанное выше позволяет сформулировать следующую проблему минимизации
суперпозиций. Пусть $H=\{\varphi_i\}$, $i\in I_0$ ($I_0$ "--- некоторое
множество индексов) "--- система отображений, а $\beta_i$ "---
<<стоимость>> отображения $\varphi_i$. (Это могут быть стоимость
изготовления блока, соответствующего $\varphi_i$, его быстродействие,
накопление погрешностей при выполнении $\varphi_i$ и~др.) Всякой
$H$-реализуемой функции $f$ можно поставить в соответствие некоторое
множество формул $\{f_j\}$, $j\in I_1$ ($I_1$ "--- некоторое множество
индексов), тождественных $f$. Пусть $n_i^j$ "--- количество блоков
$\varphi_i$ в формуле $f_j$. Тогда <<стоимость>> $c_j$ формулы $f_j$
определяется так:
\begin{equation}\tag{1.19}
c_j=\sum_{i\in I_0}n_i^j\beta_i .
\end{equation}
Задача состоит в том, чтобы найти в множестве $I_1$ такой индекс $j=j_0$,
которому соответствовало бы наименьшее значение $c_j$, т.\,е. наименьшая
стоимость определяется по формуле
\begin{equation}\tag{1.20}
c=\inf_{j\in I_1}c_j .
\end{equation}

\begin{figure}[h]
\centering
\begin{tikzpicture}[x=1cm,y=1cm]
 \coordinate (A) at (0,0); \coordinate (C) at (3.6,0);
 \coordinate (B) at (1.8,0.95); \coordinate (D) at (1.8,-0.95);
 \draw[thick] (-0.9,0) node[above]{$x$} -- (A) node[above]{$A$};
 \draw[thick] (A) -- (B) -- (C) -- (D) -- (A);
 \node[blk,fill=white,minimum size=6mm] at (B) {$f_1$};
 \node[blk,fill=white,minimum size=6mm] at (D) {$f_2$};
 \node[below=3.5mm] at (B) {$B$};
 \node[above=3.5mm] at (D) {$D$};
 \node[above] at (C) {$C$};
 \draw[wire,thick] (C) -- (4.5,0) node[above]{$y$};
\end{tikzpicture}
\figcap{Рис. 13.}
\end{figure}

Сформулированная проблема в общем случае трудно разрешима. Для частных
ситуаций она решается на основе накопленного опыта в проведении
тождественных преобразований. (Вспомним, например, школьные задачи по
алгебре и тригонометрии, которые начинаются словами <<упростить
выражение\dots>>). Некоторые результаты по минимизации булевых функций
будут приведены в \S\,10 настоящей главы.

Возможность представить одно и то же отображение $f$ различными формулами
$f_1$ и $f_2$ можно условно изобразить в виде диаграммы (рис.~13). Один и
тот же образ $y$ получится при движении входной информации (т.\,е. $x$) по
маршрутам $ABC$ и $ADC$. Такого вида диаграммы принято называть
\emph{коммутативными}~[62].

Отношение $\beta$ между отображениями $f_1$ и $f_2$ (обозначим его
$f_1\underset{\beta}{\sim}f_2$), задаваемое коммутативной диаграммой,
является отношением эквивалентности~[43]. Действительно, оно обладает
следующими свойствами:

\noindent рефлексивностью "--- $f\underset{\beta}{\sim}f$ $\forall f\in\M(H)$;\\
симметричностью "--- если $f_1\underset{\beta}{\sim}f_2$, то
$f_2\underset{\beta}{\sim}f_1$;\\
транзитивностью "--- если $f_1\underset{\beta}{\sim}f_2$,
$f_2\underset{\beta}{\sim}f_3$, то $f_1\underset{\beta}{\sim}f_3$.

Таким образом, множество $\M(H)$ $H$-реализуемых функций может быть
разбито на классы эквивалентности по признаку эквивалентности
коммутативных диаграмм.

\subsection*{Упражнения и задачи}
1. Доказать, что множество $C^m(\R^n)$ ($\forall n\in\N$) $m$ раз везде
дифференцируемых функций (всевозможного
\bookpage{33}%
конечного числа аргументов) замкнуто. Верно ли это утверждение и при
$m=\infty$?

2. Доказать, что множество $C^A(\R^n)$ ($\forall n\in\N$) везде
аналитических функций (т.\,е. раскладывающихся в ряд Тейлора с ненулевым
радиусом сходимости в любой точке $\R^n$) замкнуто.

3. Показать, что замкнутым множеством является множество $\varepsilon_n$
$n$ раз непрерывно дифференцируемых элементарных функций.

4. Построить полную систему в $\varepsilon_n$ (пример
3),\notefn{по-видимому, имеется в виду упражнение~3, в котором определено
множество $\varepsilon_n$.} содержащую конечное число функций (не считая
констант из $\R$).

5. Пусть $\varphi_i\in C(\R^{n_i})$ ($i=1,\dots,s$), где $C(\R^k)$ "---
множество всех непрерывных в $\R^k$ функций. Пусть $f$ есть суперпозиция
функций $\varphi_i$, имеющая структуру
$f=\chi[\varphi_1(x_{11},\dots,x_{n_11}),\dots,\varphi_s(x_{1s},\dots,x_{n_ss})]$.
Очевидно, что $f$ как суперпозиция непрерывных функций есть непрерывная
функция вида $f(x_1,\dots,x_n)$, где $x_1,\dots,x_n$ получается из строки
аргументов $x_{11},\dots,x_{n_11},\dots,x_{1s},\dots,x_{n_ss}$ путем
вычеркивания повторно встречающихся букв. Таким образом, приходим
к отображению $\chi^*\colon\R^{n_1}\times\dots\times\R^{n_s}\to\R^n$ (где
$\chi^*(\R^{n_1}\times\dots\times\R^{n_s})\subset C(\R^n)$ "--- образ этого
отображения), вполне определяемому структурой
$\chi$.\notefn{по смыслу отображение $\chi^*$ сопоставляет набору функций
$(\varphi_1,\dots,\varphi_s)$ функцию $f$, т.\,е. действует из
$C(\R^{n_1})\times\dots\times C(\R^{n_s})$ в $C(\R^n)$; так же следует
понимать и запись $\chi(\R^{n_1}\times\dots\times\R^{n_s})$ в следующем
абзаце "--- как образ этого отображения.}

Для каких структур $\chi(\R^{n_1}\times\dots\times\R^{n_s})=C(\R^n)$?
Привести пример структуры, для которой такое равенство не достижимо.
Выяснить, в частности, всякая ли непрерывная функция $f(x,y,z,t)$ может быть
представлена в виде $f(x,y,z,t)\equiv\varphi(x,\varphi(y,\varphi(z,t)))$,
где $\varphi(x,y)\in C(\R^2)$? В упоминавшейся в \S\,4 (пример~2)
13~проблеме Гильберта речь шла о представлении непрерывных функций трех и
большего числа аргументов в виде суперпозиций непрерывных функций двух
аргументов; при этом разрешалось использовать формулы любой структуры.
Здесь же ставится вопрос так: нельзя ли сократить множество структур,
используемых для представления непрерывных функций трех и более аргументов
через непрерывные функции двух аргументов? Может быть существует такая
одна структура, которой будет достаточно для этой цели?

6. Вопросы, поставленные в задаче~5, могут быть перенесены и на другие
функциональные классы ($C^m(\R^n)$, $H^s(\R^n)$ и~др.).
